% Speciation as catastrophe --- the multiplicative-dissipation automaton.
% Source: multiplicative-dissipation-ca-v2.tex, reduced to the model and its
% consequences.  Cut entirely: the colour rules and the hue thermostat SDE,
% the interface, the options-at-a-glance table, the numerical notes beyond two
% clauses, the fixed-constants table, and the display machinery (tags, ghosts,
% species perimeter).  Four of the source's six figures are kept.
% Renames, all forced by \S7.1's rule against contradicting a thesis-wide
% symbol, and all chosen to expose the identity with \cref{var:sec:rrrr}:
% energy \alpha -> V_i (it is that section's potential; \alpha already carries
% two sanctioned meanings), seed energy \alpha_0 -> V_0, per-cell supply
% E -> \varphi (it is that section's inflow, charged per cell), cost per
% sub-species \beta -> c (it is that section's consumption, charged per type;
% \beta is the speciation rate and nothing else), rare-type advantage
% \lambda -> \nu (\lambda is the birth rate thesis-wide), lattice side N -> L
% (N already serves as an abundance and as an integer capacity).
% Retained without change: \delta, X_i, N_i, S_t, \theta, \kappa, m, d.

\section{Speciation as catastrophe}
\label{var:sec:ca}

\subsection{The same hazard, with the inflow made endogenous}
\label{var:sec:camotivation}

\Cref{var:sec:rrrr} drove catastrophe by a potential replenished from outside
the population. This section keeps that hazard exactly and changes where the
potential comes from: it is now produced by the population itself, in
proportion to the ground it holds. Three consequences follow, and they are
the reason the model is worth the space.

The first is that the cycle stops being a cycle. When inflow rises with area
and cost rises with internal diversity, and nothing ever lowers the
diversity but competition, a population is driven onto the locus where the
two cancel and parks there. It does not oscillate; it sits at a marginal
state and waits for the hazard.

The second is that the catastrophe divides rather than destroys. The units
whose multiplicity drives the hazard are not individuals but \emph{types},
and when the hazard fires they are not lost --- they are promoted, all at
once, to independence. Nothing dies. What changes is the law by which the
survivors interact.

The third follows from the second. Because each catastrophe leaves behind
products that obey the same law, the cycle branches instead of repeating,
and no revival term is needed: at the settings used here a single founder
gives rise to some $125{,}000$ species by $t=3000$, of which about fifty are
then alive. \Cref{var:sec:rrrr} needed $\omega$ to keep
$0$ from absorbing; here the catastrophe restarts the process itself.

Set against \cref{var:sec:speciation}, the model is the other answer to the
same question. There, speciation is an elementary birth at rate $\beta(t)$,
the integrated rate converges, and $S_\infty$ has a proper limit law. Here it
is a catastrophe at rate $\delta X_i/V_i$, nothing converges, and the species
count is the statistics of a branching cascade. The two are not rival
descriptions of one phenomenon so much as the two ends of what a speciation
rate can be made to mean.

The mechanism that charges the cost per type is what
\cref{var:sec:pimentelplus} proposed and never ran. There, fitness carries a
penalty $-\gamma S^{(a)}_t$ proportional to the species a population has
accumulated \cref{var:eq:fitpenalty}; here the potential carries a penalty
$-c X_i$ proportional to the sub-species, and the model has been simulated.
What is missing here is the competitive reversal, and what is missing there
is a run.

\subsection{State}
\label{var:sec:castate}

The lattice is $L\times L$, with $L=256$ throughout, so $L^2$ cells; the
neighbourhood is Moore and the boundary periodic, so the lattice is a torus
with no edge and no corner. Each cell is empty or occupied.

Every occupied cell carries two nested labels, a \textbf{species} and, within
it, a \textbf{sub-species}. Only the sub-species is stored, the species being
recovered from it, so that when a species fragments its sub-species are
re-pointed at their new parents and every cell follows without being touched.
For each extant species $i$ the model carries
\begin{center}
\begin{tabular}{@{}ll@{}}
$V_i$ & accumulated potential, a real variable obeying \cref{var:eq:caode} \\
$N_i$ & cells occupied by species $i$ \\
$X_i$ & number of \emph{extant} sub-species of species $i$ \\
\end{tabular}
\end{center}
with $X_i\le N_i$, equality holding when every cell is its own sub-species.
Both $N_i$ and $X_i$ are read off the lattice.

That $X_i$ counts what is alive rather than what has ever existed is not a
bookkeeping detail. A sub-species driven extinct by competition inside its
parent stops costing that parent anything, so the coarsening of
\cref{var:sec:cacompetition} is a genuine regulator of dissipation and a
species can hold its diversity, and hence its potential, in balance. Under a
cumulative count the cost could only ever rise, and fragmentation would
become certain from the moment a species began to diversify.

The run starts from one occupied cell at the centre, of one species with one
sub-species, at potential $V_0$. Time is counted in \textbf{sweeps}: one
elementary contest advances the clock by $L^{-2}$, so a sweep is one attempt
per site on average. Competition is a discrete sequence of events while
mutation and the potential are continuous, and the two are separated by
advancing in chunks of a few hundred to a few thousand contests and applying
the continuous updates once over each. That splitting is the model's only
approximation in time, and its error is bounded by the movement of $V_i$
within one chunk --- a few parts in a thousand at the settings used here.

\subsection{Competition}
\label{var:sec:cacompetition}

At each iteration a cell is drawn uniformly from all $L^2$ and one of its
eight neighbours uniformly from the eight. If the two carry the same label,
or both are empty, nothing happens. Otherwise they compete, and the winner's
label is copied into the loser's cell.

Contests between occupied cells are one rule. Give each occupied cell $\zeta$,
carrying sub-species $u$ of species $i$, the strength
\begin{equation}
w(\zeta)=\underbrace{\vphantom{(n_\zeta+1)^\theta}V_i}_{\text{potential}}
\times\underbrace{(n_\zeta+1)^{\theta}}_{\text{local support}}
\times\underbrace{\ee^{-\nu|u|/\kappa_i}}_{\text{rare-type discount}},
\qquad
\kappa_i=\frac{N_i}{X_i},
\label{var:eq:castrength}
\end{equation}
where $n_\zeta\in\{0,\dots,8\}$ counts the neighbours matching $\zeta$ at the
level of the contest, $|u|$ is the number of cells sub-species $u$ occupies,
and $\kappa_i$ is the mean sub-species area within the parent. Let the
contest go to $\zeta$ with probability $w(\zeta)/(w(\zeta)+w(\zeta'))$. Two
readings follow.

\begin{itemize}[leftmargin=1.6em,itemsep=0.3em,topsep=0.3em]
\item \textbf{Against a sub-species of the same species.} The parent's
potential cancels from numerator and denominator, as it must, since two
sub-species share one store. What survives is support and rarity, and
$n_\zeta$ counts neighbours matching at the sub-species level.
\item \textbf{Against a distinct species.} The rarity factor is omitted and
$n_\zeta$ counts neighbours of the same \emph{species}, so that for species
$A$ and $B$
\begin{equation}
\pr{A\text{ wins}}=\frac{V_A(n_A+1)^{\theta}}
{V_A(n_A+1)^{\theta}+V_B(n_B+1)^{\theta}} .
\label{var:eq:catypeIII}
\end{equation}
This is the sense in which $V_i$ is a potential: it buys territory, and it
buys it only against other species.
\end{itemize}

Rarity is withheld between species deliberately. Inside a species it
preserves the diversity that drives the model; between species it would
merely blunt $V_i$, which is the quantity the model is about.

Contests against empty ground stand outside \cref{var:eq:castrength}
entirely. Empty ground is a rival of fixed strength $d$, so an occupied cell
takes a contested empty neighbour with probability $1-d$ and otherwise
nothing happens --- the occupied cell is never replaced, and empty ground
never invades. Colonisation carries no dependence on $V_i$, which is a
deliberate restriction: a species cannot widen its frontier by being rich,
only take ground already held.

The support exponent $\theta$ buys surface tension. The factor
$(n_\zeta+1)^\theta$ ranges over $[1,9^\theta]$, so at $\theta=3$ a cell with
all eight neighbours of its own kind beats an isolated intruder better than
$998$ times in $1000$: a domain pays along its perimeter and nothing in its
interior, so it rounds up and can be eaten only from the edge. Without it
boundaries wander and domains dissolve into filigree
(\cref{var:fig:ca-support}).

The rarity exponent $\nu$ decides whether a burst is a catastrophe or a
renaming, which is a sharper role than it looks. Since $|u|$ enters
negatively, the \emph{larger} sub-species is penalised, and at $\nu=1.14$ one
exceeding its rival by a mean sub-species area wins only about a quarter of
their contests. What this changes is not how many sub-species a species
carries but how large any one may grow. Set $\nu=0$ and coarsening inside a
species does what a voter model does --- one sub-species engulfs the domain,
and of the $7262$ fragments the largest species would yield at $t=2000$
exactly one is big enough to survive. At $\nu=1.14$ the domain is instead a
mosaic of comparable patches, and about half of $1534$ fragments clear the
threshold. Diversity in the bare sense survives either way: $X_i$ climbs to
the balance point under both. Only under the advantage does fragmentation
actually multiply the species.

\begin{figure}[!htbp]
\centering
\begin{minipage}[t]{0.46\textwidth}\centering
\includegraphics[width=\linewidth]{fig08_ca/support-on}\\[3pt]
{\small support on, $\theta=3$}
\end{minipage}\hfill
\begin{minipage}[t]{0.46\textwidth}\centering
\includegraphics[width=\linewidth]{fig08_ca/support-off}\\[3pt]
{\small support off, $\theta=0$}
\end{minipage}
\caption{The same seed at the same time ($t=2000$), with neighbour support on
and off. With it, domains are compact and their boundaries short: seven per
cent of occupied cells lie on a species boundary. Without it, boundaries
wander, domains interpenetrate, and the boundary fraction rises to thirty-six
per cent. The ragged state also carries far more species at once --- $1335$
against $26$ --- because nothing gathers a domain together and small holdings
are never squeezed out. Hue labels the taxonomy and nothing else: it is
carried along by species the competition rules have already made, and no cell
ever changes hands because of a colour.}
\label{var:fig:ca-support}
\end{figure}

\subsection{Mutation and the potential}
\label{var:sec:camutation}

A sub-species arises when one cell mutates inside its species' domain,
keeping its species label and taking a new sub-species label, so that
$X_i\mapsto X_i+1$; thereafter it spreads by the within-species contest. The
rate is $m V_i$ per cell, giving a species-wide rate
\begin{equation}
R^{\text{mut}}_i=m\,V_i N_i ,
\label{var:eq:camut}
\end{equation}
which couples novelty to prosperity: rich species diversify faster and
thereby hasten their own dissipation, which is the feedback the model exists
to display. The uncoupled alternative $mN_i$ is available and changes the
model's character, not merely its rate. One consequence deserves note: if the
mutating cell was the last of its sub-species, that sub-species dies as the
new one is born and $X_i$ is unchanged. Mutation adds diversity only where
there is something to mutate away from.

Each occupied cell supplies potential at rate $\varphi$ and each extant
sub-species costs $c$ --- the cost of multiplicity, charged per type rather
than per head. Hence
\begin{equation}
\ddt{V_i}=\varphi N_i-c\,X_i ,
\label{var:eq:caode}
\end{equation}
floored at zero, against \cref{var:sec:rrrr}'s
$\dd V/\dd t=\phi-cX_t$. The two differences are the whole model: the inflow
is charged per cell rather than supplied from outside, and the outflow is
charged per type rather than per individual.

Three readings of \cref{var:eq:caode} are worth recording.
\begin{enumerate}[label=(\roman*),leftmargin=2.1em,itemsep=0.25em,topsep=0.25em]
\item $V_i$ does not appear on the right. It is a pure accumulator with no
fixed point of its own, growing while $X_i<\varphi N_i/c$ and falling once
diversity passes that threshold.
\item Since $X_i\le N_i$, the most negative attainable drift is
$N_i(\varphi-c)$, attained at $X_i=N_i$. A species can therefore lose
potential at all only if $c>\varphi$, and that is the condition under which
the catastrophe of \cref{var:sec:cafragment} is reachable. Below it there is
no model: nothing ever bursts, and the run reduces to the lattice filling and
then relabelling indefinitely.
\item A species carrying a single sub-species needs at least $c/\varphi$
cells to break even --- about nine at the settings used here. Below that area
its potential drains away whatever it does, which is what governs the fate of
the smallest fragments of a burst.
\end{enumerate}

A species at $V_i=0$ is not dead. It loses every between-species contest it
enters against anyone with potential to spend, and stops mutating, but it
still takes empty ground, and if its area exceeds $cX_i/\varphi$ its
potential recovers. It is merely defenceless until it does.

\subsection{Fragmentation}
\label{var:sec:cafragment}

When a species has come to carry too many sub-species it \textbf{fragments}:
every sub-species breaks away at once and becomes an independent species. The
domains do not move. What changes is the law by which they interact, which
switches from the within-species contest to \cref{var:eq:catypeIII}:
neighbouring patches that a moment earlier were contesting a shared
inheritance now contest each other in proportion to their separate
potentials.

Two triggers are live at once and tested in order. A species fragments
deterministically if the update to \cref{var:eq:caode} would take $V_i$ to
zero or below; failing that it fragments with hazard
\begin{equation}
R^{\text{frag}}_i=\delta\,\frac{X_i}{V_i},
\label{var:eq:cafrag}
\end{equation}
which is \cref{var:eq:rrrr}'s catastrophe rate $\delta X_t/V(t)$ with the
count of individuals replaced by the count of types. Because
\cref{var:eq:cafrag} diverges as $V_i\downarrow0$ it anticipates the
deterministic trigger rather than competing with it, while also permitting
early bursts in a species that is merely strained. At the settings used here
the hazard does essentially all the work: a species in the regime of
\cref{var:sec:caconsequences} holds a potential of several hundred and never
approaches zero, so it bursts because \cref{var:eq:cafrag} fires and not
because it is bankrupt. A species carrying one sub-species cannot fragment,
there being nothing for a division to mean; it sits at $V_i=0$ until its area
recovers or it is overrun.

Each new species takes exactly the cells its sub-species held and begins with
$X=1$, its own only sub-species. Its potential is the one thing the burst
does not settle by itself, the parent's store being wound up rather than
divided by any natural rule, and it is reset to $V_0$. Two splitting rules
were also implemented --- equal shares, and shares proportional to area ---
but both interact awkwardly with the deterministic trigger, which fires
precisely when there is nothing left to divide and so hands every fragment a
potential of zero. Resetting is the only rule that behaves identically under
both triggers.

\begin{figure}[!htbp]
\centering
\begin{minipage}[t]{0.46\textwidth}\centering
\includegraphics[width=\linewidth]{fig08_ca/burst-before}\\[3pt]
{\small before}
\end{minipage}\hfill
\begin{minipage}[t]{0.46\textwidth}\centering
\includegraphics[width=\linewidth]{fig08_ca/burst-after}\\[3pt]
{\small after}
\end{minipage}
\caption{A burst, across the single frame in which it happens, at
$t\approx1239$. The species held $12{,}826$ cells --- a fifth of the lattice
--- and still had $V_i=755$ in hand when \cref{var:eq:cafrag} fired; all
$1444$ of its sub-species became species at once, and it was the only species
to fragment in that frame. Nothing has moved: the outline of the parent is
exactly preserved and its neighbours are untouched. What has changed is the
law inside that outline. Each fragment is drawn in a fresh hue here so that
the division is unmistakable.}
\label{var:fig:ca-burst}
\end{figure}

\subsection{What the simulations show}
\label{var:sec:caconsequences}

\Cref{var:fig:ca-sequence,var:fig:ca-founder} are exact realisations at the
settings above. Four consequences follow from the rules as stated.

\paragraph*{The occupied region never retreats.} No rule vacates a cell:
contests against empty ground cannot, by construction, and contests between
occupied cells only relabel. The occupied set is monotonically
non-decreasing, and the lattice fills --- here by about $t=1100$. After that
the dynamics reduce to relabelling, a biased voter model on a full lattice
punctuated by mutation and fragmentation.

\paragraph*{The population sits at the balance point.} The cycle does not run
away. Because mutation raises $X_i$ while nothing lowers $\varphi N_i$, a
species is driven towards
\begin{equation}
X_i\approx\frac{\varphi N_i}{c},
\label{var:eq:cabalance}
\end{equation}
at which production and dissipation cancel and $V_i$ goes flat.
\Cref{var:fig:ca-founder} is that statement in one picture: the founder's
potential climbs steeply while diversity is still cheap, flattens near $730$
once $X_i$ reaches \cref{var:eq:cabalance} at about $t=175$, and stays there.
Every substantial species in a mature run is found on that locus to within a
few per cent, with its potential parked at a few hundred and its drift near
zero. The system settles into a marginal state rather than an oscillation,
and what ends a species is not bankruptcy but the accumulating hazard of
\cref{var:eq:cafrag}, which grows with the diversity the balance point
obliges it to carry.

\paragraph*{Fragments are filtered by area.} A burst is not a fair division.
Each fragment inherits its sub-species' cells, and those areas are broadly
distributed; but by reading (iii) of \cref{var:sec:camutation} a fragment of
fewer than $c/\varphi$ cells cannot pay for even its own single sub-species.
Its potential drains within a fraction of a sweep, after which it loses every
contest it enters and is overrun. Immediately after a large burst, hundreds
of species may sit at $V_i=0$ at once. The great majority of what a burst
creates is gone almost immediately, and only the larger fragments survive to
grow, diversify and burst in their turn --- which is why a run that has
produced tens of thousands of species holds only tens alive.

\paragraph*{Growth and fragmentation overlap.} At this colonisation rate the
first turn of the cycle completes long before the lattice is full: the
founder of \cref{var:fig:ca-founder} bursts on $8689$ cells, a seventh of the
ground, with open country still ahead of it. A species reaches the diversity
that destroys it before it reaches the walls, so the run is never a sequence
of clean generations.

\begin{figure}[!htbp]
\centering
\begin{minipage}[t]{0.235\textwidth}\centering
\includegraphics[width=\linewidth]{fig08_ca/seq-250}\\[3pt]{\small $t=250$}
\end{minipage}\hfill
\begin{minipage}[t]{0.235\textwidth}\centering
\includegraphics[width=\linewidth]{fig08_ca/seq-600}\\[3pt]{\small $t=600$}
\end{minipage}\hfill
\begin{minipage}[t]{0.235\textwidth}\centering
\includegraphics[width=\linewidth]{fig08_ca/seq-1100}\\[3pt]{\small $t=1100$}
\end{minipage}\hfill
\begin{minipage}[t]{0.235\textwidth}\centering
\includegraphics[width=\linewidth]{fig08_ca/seq-3000}\\[3pt]{\small $t=3000$}
\end{minipage}
\caption{One run on a $256^2$ lattice. At $t=250$ a single species holds six
per cent of the ground and already carries $443$ sub-species, visible as
mottling within one narrow band of hue. By $t=600$ it has burst six times and
the fragments are still spreading into open ground. The lattice closes at
about $t=1100$; by $t=3000$ the run holds some fifty living species out of
$125{,}000$ ever created.}
\label{var:fig:ca-sequence}
\end{figure}

\begin{figure}[!htbp]
\centering
\includegraphics[width=0.82\textwidth]{fig08_ca/fig08}
\caption{The founder species, followed from its first cell to the burst that
ends it at $t=354$. Potential $V_i$ above, sub-species count $X_i$ below
against the balance point $\varphi N_i/c$ of \cref{var:eq:cabalance}; the
dotted line marks the burst. The lower panel is the argument of
\cref{var:sec:caconsequences} in one picture --- the dashed curve is area,
rescaled, and diversity tracks it. The species is not bankrupt when it
bursts: it still holds $V_i\approx730$ on $8689$ cells, but by then it
carries $964$ sub-species and the hazard \cref{var:eq:cafrag} has grown large
enough to fire.}
\label{var:fig:ca-founder}
\end{figure}

\paragraph{Open questions.}\label{var:rem:open-ca}
\begin{enumerate}[leftmargin=1.6em,itemsep=0.2em,topsep=0.2em]
\item \textbf{Whether the balance point restores tractability.} This is the
sharpest question the model raises. On \cref{var:eq:cabalance} the potential
is flat, so $V_i$ is slowly varying at some $\bar V$, and the hazard
\cref{var:eq:cafrag} becomes
$\delta X_i/V_i\approx(\delta\varphi/c\bar V)N_i$ --- proportional to the
species' \emph{area}. But a catastrophe rate proportional to size is exactly
\ChCore{}'s $\delta X_t$. Whether the species count can therefore be read as
a birth--death--catastrophe process with burst size $X_i$, with an effective
$\delta$ of $\delta\varphi/c\bar V$, is the question to go after: it would
put the chapter's least tractable model back onto the chapter's own
machinery, and it is a conjecture supported by
\cref{var:fig:ca-founder} rather than anything derived.
\item \textbf{The statistics of the cascade.} \Cref{var:sec:rrrr}'s open
question was the distribution of the interval between catastrophes. Here the
catastrophes branch, so the analogous object is the distribution of the
number of \emph{viable} fragments per burst --- those clearing the
$c/\varphi$ threshold --- and whether its mean sits above or below one
decides whether the species count grows or collapses. The run holds tens
alive out of $125{,}000$ created, which suggests a mean near one and a
process near criticality, but that has not been measured.
\item \textbf{A continuous species concept.} The model could be defined by
intervals of hue rather than by discrete labels, with species identity a
matter of proximity on the colour circle and speciation a matter of a
lineage drifting far enough from its relatives to fall outside the interval.
That would remove the taxonomy as a primitive and make it an observable, and
it is not obvious whether the balance point survives the change.
\item \textbf{What it could be fitted to.} As with \cref{var:sec:rrrr}, the
model has never been fitted, and $\varphi$, $c$ and $\delta$ enter the hazard
only through the combination $\delta\varphi/c$ once the balance point is
reached. A time series of the species count would not separate them; what
second observable would is not known.
\end{enumerate}

\FloatBarrier
